%! AP Statistics — Unit 8, Lesson 8.1 — Guided Notes
\documentclass[10pt]{article}
\usepackage{apstats-article}
\usepackage{apstats-boxes}

\begin{document}

\pageheader{Unit 8, Lesson 8.1}{Guided Notes}
\namedateperiod

\begin{objectivebox}
By the end of this lesson, I will be able to:
\begin{itemize}
    \item[$\square$] \writeline
    \item[$\square$] \writeline
    \item[$\square$] \writeline
\end{itemize}
\end{objectivebox}

\begin{vocabbox}
\termblanklong{Observed counts}
\termblanklong{Expected counts}
\termblanklong{Chi-square test}
\termblanklong{Distribution of a categorical variable}
\end{vocabbox}

\begin{hookbox}
\textbf{Discussion:} A bag of candies claims 20\% of each of 5 colors. You open a bag
and count: Red 14, Orange 8, Yellow 12, Green 9, Blue 7 (total = 50). Is your bag unusual,
or is this just normal variation? How would you decide?

\writelines{2}
\end{hookbox}

\begin{notesbox}{What Are Expected Counts?}
If we \emph{assume} a categorical variable follows a specific distribution, the
\textbf{expected count} for each category is:

\[
    \text{Expected count} = n \times p_{\text{category}}
\]

where $n$ is the \blank{3cm} and $p_{\text{category}}$ is the \blank{5cm}.

\vspace{0.1in}
\textbf{Example:} A coffee shop sells 200 drinks. If each of 4 flavors is equally popular
(25\% each), the expected count for each flavor is $200 \times 0.25 = $ \blank{2cm}.

\vspace{0.1in}
Expected counts represent what we would see \blank{6cm}
if the assumed distribution were true. They are not necessarily whole numbers.
\end{notesbox}

\begin{notesbox}{What Does ``Surprising'' Mean Statistically?}
After computing expected counts, we compare them to the \blank{3cm} counts.

\vspace{0.1in}
\textbf{The gap} for a single category $= $ \blank{4cm}.

\vspace{0.1in}
A gap of \blank{1cm} means perfect agreement. Larger gaps indicate that the data
\blank{6cm}.

\vspace{0.1in}
\renewcommand{\arraystretch}{2}
\begin{tabularx}{\linewidth}{Xcc}
    \textbf{Category} & \textbf{Observed} & \textbf{Gap (Obs $-$ Exp)} \\
    \hline
    Espresso   & 48 & \blank{2cm} \\
    Latte      & 72 & \blank{2cm} \\
    Cappuccino & 36 & \blank{2cm} \\
    Cold Brew  & 44 & \blank{2cm} \\
\end{tabularx}

\vspace{0.1in}
Which category has the largest gap? \blank{3cm} \quad
Is the gap positive or negative? \blank{3cm}
\end{notesbox}

\begin{notesbox}{Preview: The Chi-Square Idea}
Looking at one gap at a time is useful, but a chi-square test considers
\blank{4cm} categories simultaneously.

\vspace{0.1in}
The general idea: we \blank{4cm} the gaps across all categories (after adjusting
each gap for the size of the expected count). The resulting number is called the
\textbf{chi-square statistic}, denoted $\chi^2$ (``chi-square'').

\begin{itemize}
    \item A $\chi^2$ value close to \blank{2cm} means observed counts closely match
          expected counts.
    \item A \blank{2cm} $\chi^2$ value means the observed counts differ substantially
          from expected.
    \item We will learn the formula in Lesson \blank{2cm}.
\end{itemize}
\end{notesbox}

\begin{notesbox}{Unit 8 Arc Overview}
\renewcommand{\arraystretch}{1.5}
\begin{tabularx}{\linewidth}{lX}
    \textbf{Lesson 8.1} & Framing: what question does Unit 8 ask? \\
    \textbf{Lesson 8.2} & Setting up the chi-square \blank{4cm} test \\
    \textbf{Lesson 8.3} & Carrying out the chi-square GOF test \\
    \textbf{Lessons 8.4--8.5} & Chi-square test for \blank{4cm} (two-way tables, multiple groups) \\
    \textbf{Lesson 8.6} & Chi-square test for \blank{4cm} (two categorical variables, one sample) \\
    \textbf{Lesson 8.7} & Synthesis: choosing the right chi-square procedure \\
\end{tabularx}
\end{notesbox}

\begin{practicebox}
\textbf{Guided Practice:} A school surveys 120 students about their preferred study
environment. Results: Library 38, Home 52, Caf\'e 18, Classroom 12.

\begin{enumerate}
  \item If preference were equally distributed among 4 options, what is the
        expected count for each category?
        \writelines{2}

  \item Compute the gap (Observed $-$ Expected) for each category.
        \writelines{2}

  \item Without any formal calculation, does this data suggest students do
        \emph{not} have equal preferences? Explain.
        \writelines{2}
\end{enumerate}
\end{practicebox}

\end{document}
